\chapter{Command-Line Interface and API Reference}
\label{app:api}

This appendix provides a practical reference for using the TimeFlowCodec command-line interface (CLI) and Python developer API.

\section{Command-Line Interface (\texttt{tfc})}
The package installs a unified console executable \texttt{tfc}:

\subsection{Encoding Video to .tfc}
\begin{verbatim}
tfc encode input.mp4 output.tfc [OPTIONS]

Options:
  --preset [screen|anime|lownoise|general]
                        Select encoding configuration preset. Default: screen
  --tau FLOAT           Override error ratio threshold (e.g. 0.0001).
  --slope FLOAT         Override slope threshold (default: 0.001).
  --tiling INT          Set spatial block dimension (1 for per-pixel).
  --window INT          Temporal chunk duration in frames (v5 default: 48).
  --comp [none|zlib|lzma|zstd]
                        Payload compression algorithm. Default: zstd
  --max-ram-mb INT      Maximum RAM allocated before disk memmap fallback.
\end{verbatim}

\subsection{Decoding .tfc to Video}
\begin{verbatim}
tfc decode input.tfc output.mp4 [OPTIONS]

Options:
  --fps FLOAT           Override playback frame rate if unspecified in v1-v4.
  --lossless            Force lossless RGB output encoding via libx264rgb.
\end{verbatim}

\section{Python Developer API}
The high-level programmatic interface provides simple functional bindings:

\begin{verbatim}
import timeflowcodec as tfc

# Encode MP4 to .tfc container
tfc.encode_video_to_tfc(
    in_path="recording.mp4",
    out_path="recording.tfc",
    preset="screen",
    container_version=5,
    payload_comp_type=tfc.COMP_ZSTD,
)

# Decode .tfc container to uncompressed in-memory RGB array
rgb_frames = tfc.decode_tfc_to_video(
    in_path="recording.tfc",
    out_path=None  # Returning array directly in memory
)
print("Decoded video tensor shape:", rgb_frames.shape)
\end{verbatim}
